\documentclass[10pt,a4paper]{amsart}
\usepackage[margin=19mm]{geometry}
\usepackage{amsmath,amssymb,mathtools}
\usepackage[scr=rsfs,cal=euler]{mathalfa}
\usepackage{xfrac}
\usepackage[unicode,hidelinks,bookmarks=false]{hyperref}
\newcommand{\ie}{i.e.\ }
\newcommand{\cf}{cf.\ }
\newcommand{\resp}{respectively\ }
\newcommand{\Subsection}{\subsection}
\providecommand{\nobreakdash}{\nobreak\hbox{-}}
\setlength{\emergencystretch}{2em}
\multlinegap0pt
\usepackage{dsfont}
\usepackage{amssymb}
\usepackage{braket}
\newtheorem{lemma}{Lemma}
\title{Partitions of unity and barycentric algebras}
\author{Anna Zamojska-Dzienio}
\date{Source: arXiv:2501.00937v1 (CC BY 4.0)}
\begin{document}
\maketitle

\noindent\emph{Source and licence.} Partitions of unity and barycentric algebras. Version arXiv:2501.00937v1 (CC BY 4.0), distributed by arXiv under the Creative Commons Attribution 4.0 International licence, http://creativecommons.org/licenses/by/4.0/, as recorded in the arXiv licence metadata for this exact version and shown on the abstract page https://arxiv.org/abs/2501.00937v1. Full licence notice: \texttt{licenses/arxiv-2501.00937-CC-BY-4.0.txt}.\par\smallskip

\noindent\textit{Editorial scope.} Counted unit: the article's lemma lemma (source0.tex lines 317--319). The article's complete original proof of it is delivered with it (source0.tex lines 321--331, one \texttt{proof} environment of 11 lines). No mathematics is written, repaired or re-derived: the statement and the proof are the article's own lines, in the article's own notation, and no retained line is substituted or re-flowed.  No further source-local context is needed: every cross-reference of every retained line resolves inside the two delivered spans.  Nothing was omitted from any retained span, so the package carries no omission record.  No retained line contains a citation, so the wrapper carries no bibliography and no bibliography entry.  No retained line contains a figure environment, an image-inclusion command or drawing code, so no image is delivered and the package declares no asset dependency.  Editorial formatting only, in addition: an AMS \texttt{amsart} preamble that restates the article's own declarations and macros used by the retained lines, namely the environments \texttt{lemma} on separate counters, together with the standard packages those lines need.  The retained environments print in this document as Lemma 1 in order of appearance, whereas the article prints the same items with the numbering of its own full document.  The complete native source of the article as distributed in the arXiv e-print for this exact version is bundled unchanged as \texttt{sources/171035/source0.tex}.
\begin{lemma}
$\mathbf{Set}^{\mathds 1}(\Pi,\mathcal{I}^n)$ is a subalgebra of the barycentric algebra $\mathbf{Set}(\Pi,\mathcal{I}^n)$.
\end{lemma}

\begin{proof}
For $f,g\in\mathbf{Set}^{\mathds 1}(\Pi,\mathcal{I}^n)$ and $q\in I^\circ$, one has
\begin{align*}
\underline{q}(f,g)&=\underline{q}\left((f_1,f_2,\ldots ,f_n),(g_1,g_2,\ldots g_n)\right)\\
&=\left(\underline{q}(f_1,g_1),\underline{q}(f_2,g_2),\ldots \underline{q}(f_n,g_n)\right).
\end{align*}
Then
\begin{align*}
\sum_{i=1}^{n}\underline{q}(f_i,g_i)=\underline{q}\left(\sum_{i=1}^{n}f_i,\sum_{i=1}^{n}g_i\right)=\underline{q}(\mathds 1_\Pi,\mathds 1_\Pi)=\mathds 1_\Pi.
\end{align*}
\end{proof}

\end{document}
